\section{Week 3: Fitting Models by Least Squares: Error Surfaces and
Optimization}

\begin{quote}
\textbf{Chapter 2: Parameter estimation methods} · Lecture 1 h · Lab 2 h
\textbf{This week's running data:} the simulated forgetting experiment
from Week 2 (power model, \(A^\ast = 0.9\), \(d^\ast = 0.5\), \(n = 50\)
items at 8 intervals). We refit it properly and compare with last week's
linearized fits.
\end{quote}

\begin{center}\rule{0.5\linewidth}{0.5pt}\end{center}

\subsection{Key terms}

Important terms of this lesson, with their Persian equivalents (repeated
even if they appeared in earlier weeks).

\begin{longtable}[]{@{}
  >{\raggedright\arraybackslash}p{(\linewidth - 4\tabcolsep) * \real{0.3333}}
  >{\raggedright\arraybackslash}p{(\linewidth - 4\tabcolsep) * \real{0.3333}}
  >{\raggedright\arraybackslash}p{(\linewidth - 4\tabcolsep) * \real{0.3333}}@{}}
\caption{\centering Key terms}\tabularnewline
\toprule\noalign{}
\begin{minipage}[b]{\linewidth}\raggedright
English
\end{minipage} & \begin{minipage}[b]{\linewidth}\raggedright
\RL{\CMfa فارسی}
\end{minipage} & \begin{minipage}[b]{\linewidth}\raggedright
\RL{\CMfa توضیح کوتاه}
\end{minipage} \\
\midrule\noalign{}
\endfirsthead
\toprule\noalign{}
\begin{minipage}[b]{\linewidth}\raggedright
English
\end{minipage} & \begin{minipage}[b]{\linewidth}\raggedright
\RL{\CMfa فارسی}
\end{minipage} & \begin{minipage}[b]{\linewidth}\raggedright
\RL{\CMfa توضیح کوتاه}
\end{minipage} \\
\midrule\noalign{}
\endhead
\bottomrule\noalign{}
\endlastfoot
fit / fitting & \RL{\CMfa برازش} & \RL{\CMfa تنظیم پارامترها تا مدل با داده جور
شود} \\
least squares & \RL{\CMfa کمترین مربعات} & \RL{\CMfa روش برازش با کمینه کردن مجموع
مربع خطاها} \\
error surface & \RL{\CMfa سطح خطا} & \RL{\CMfa خطا به‌عنوان تابعی از پارامترها} \\
optimizer & \RL{\CMfa بهینه‌ساز} & \RL{\CMfa الگوریتمی که کمینه‌ی تابع را پیدا
می‌کند} \\
parameter & \RL{\CMfa پارامتر} & \RL{\CMfa کمیت قابل تنظیم مدل که بین افراد یا شرایط
فرق می‌کند} \\
estimate / estimator & \RL{\CMfa برآورد} / \RL{\CMfa برآوردگر} & \RL{\CMfa مقدار به‌دست‌آمده
برای پارامتر} / \RL{\CMfa روش به دست آوردن آن} \\
forgetting curve & \RL{\CMfa منحنی فراموشی} & \RL{\CMfa کاهش یادآوری با گذشت
زمان} \\
exponential / power function & \RL{\CMfa تابع نمایی} / \RL{\CMfa توانی} & \RL{\CMfa دو شکل
رایج برای منحنی فراموشی و یادگیری} \\
SSE / RMSD & \RL{\CMfa مجموع مربعات خطا} / \RL{\CMfa ریشه‌ی میانگین مربع انحراف} &
\RL{\CMfa معیارهای فاصله‌ی داده و پیش‌بینی} \\
noise & \RL{\CMfa نویز} & \RL{\CMfa تغییرات تصادفی در شواهد یا در داده} \\
parameter trade-off & \RL{\CMfa بده‌بستان پارامترها} & \RL{\CMfa تغییر یک پارامتر که
با پارامتر دیگر جبران می‌شود} \\
grid approximation & \RL{\CMfa تقریب شبکه‌ای} (grid) & \RL{\CMfa محاسبه‌ی پسین روی
شبکه‌ای از مقادیر} \\
Nelder--Mead & \RL{\CMfa روش} Nelder--Mead & \RL{\CMfa بهینه‌ساز بدون نیاز به
مشتق} \\
Bayes factor (BF) & \RL{\CMfa عامل بیز} (Bayes factor) & \RL{\CMfa نسبت درست‌نمایی‌های
حاشیه‌ای دو مدل} \\
local minimum & \RL{\CMfa کمینه‌ی محلی} & \RL{\CMfa نقطه‌ای پایین‌تر از اطراف ولی نه
بهترین نقطه} \\
bounds / constraints & \RL{\CMfa کران‌ها} / \RL{\CMfa قیدها} & \RL{\CMfa محدودیت‌های مجاز
برای مقدار پارامترها} \\
reparameterization & \RL{\CMfa بازپارامتری‌سازی} & \RL{\CMfa نوشتن پارامتر در مقیاس
دیگر، مثلاً} \ensuremath{\lambda} = e\^{}\ensuremath{\eta} \\
multi-start & \RL{\CMfa شروع چندگانه} (multi-start) & \RL{\CMfa اجرای بهینه‌ساز از
چند نقطه‌ی شروع} \\
maximum likelihood estimate (MLE) & \RL{\CMfa برآورد بیشینه‌ی درست‌نمایی} &
\RL{\CMfa پارامترهایی که داده را محتمل‌ترین می‌کنند} \\
likelihood & \RL{\CMfa درست‌نمایی} & \RL{\CMfa احتمال داده به‌عنوان تابعی از
پارامترها} \\
gradient & \RL{\CMfa گرادیان} & \RL{\CMfa جهت و شیب تغییر تابع} \\
convergence & \RL{\CMfa همگرایی} & \RL{\CMfa رسیدن زنجیره‌ها به توزیع هدف} \\
large language model (LLM) & \RL{\CMfa مدل بزرگ زبانی} & \RL{\CMfa شبکه‌ی عصبی‌ای که
کلمه‌ی بعدی متن را پیش‌بینی می‌کند} \\
scaling law & \RL{\CMfa قانون مقیاس} & \RL{\CMfa رابطه‌ی توانی بین اندازه‌ی مدل و
خطا} \\
\end{longtable}

\begin{center}\rule{0.5\linewidth}{0.5pt}\end{center}

\subsection{Learning outcomes}

By the end of this week, students will be able to:

\begin{enumerate}
\def\labelenumi{\arabic{enumi}.}
\tightlist
\item
  \textbf{Define} SSE, MSE and RMSD as discrepancy functions and
  \textbf{explain} what minimizing them assumes about the noise.
\item
  \textbf{Visualize} an error surface for a two-parameter model and
  \textbf{interpret} its shape (minimum, valleys, plateaus) in terms of
  parameter trade-offs.
\item
  \textbf{Diagnose} and \textbf{remedy} common fitting failures: bad
  starting values, plateaus, local minima, and out-of-range parameters
  (using bounds, reparameterization and multi-start).
\item
  \textbf{Evaluate} a least-squares fit with RMSD and residuals, and
  \textbf{interpret} the fitted parameters.
\end{enumerate}

\begin{center}\rule{0.5\linewidth}{0.5pt}\end{center}

\subsection{Motivation \& big picture}

Last week we fitted forgetting curves by drawing a straight line through
log-transformed proportions. The power model was the true model, yet its
estimates were noticeably off (\(\hat A = 0.80\), \(\hat d = 0.37\)
instead of \(0.9\) and \(0.5\)). The method also breaks whenever a count
is zero, and it minimizes errors on the wrong (log) scale.

The fix is conceptually simple: \textbf{choose the parameters that make
the model's predictions as close as possible to the data, on the data's
own scale.} ``As close as possible'' needs a number, a \emph{discrepancy
function}. ``Choose the parameters'' needs a search procedure, an
\emph{optimization algorithm}. This week covers both.

There is a deeper reason this matters for cognitive science. Model
parameters are our \emph{measurements} of latent constructs (decay rate,
caution, sensitivity). A poor fitting procedure gives poor measurements,
and every conclusion drawn from them (group differences, correlations
with brain activity, model comparisons) inherits that error. Worse, as
we will see, an optimizer can report ``success'' while returning
nonsense.

\paragraph{Where is this used?}

\begin{itemize}
\tightlist
\item
  \textbf{Every model-based study in psychology} has a fitting step.
  Learning curves, forgetting curves, psychometric functions (how
  detection improves with stimulus strength) and dose--response curves
  in psychopharmacology are all fitted by minimizing a discrepancy.
\item
  \textbf{Clinical and developmental research} compares fitted
  parameters between groups (e.g.~the forgetting rate of patients
  vs.~controls). If the fit is poor or stuck at a bound, the group
  comparison is wrong.
\item
  \textbf{Neuroscience} fits tuning curves of neurons and hemodynamic
  response functions in fMRI with exactly these optimizers.
\item
  \textbf{Machine learning} trains models by minimizing a loss with
  gradient methods; the vocabulary (loss surface, learning rate, local
  minima) is the same as this week's.
\end{itemize}

\begin{center}\rule{0.5\linewidth}{0.5pt}\end{center}

\Needspace{11\baselineskip}

\subsection{Lecture notes}

\paragraph{In plain words}

A model has knobs (parameters). Fitting means turning the knobs until
the model's curve passes as close as possible to the data points. To do
this with a computer we need two things. First, a single number that
says how bad the current fit is: the \textbf{sum of squared errors
(SSE)}, the squared vertical distances between data points and the
curve, added up. Second, a strategy to turn the knobs so that this
number goes down: an \textbf{optimizer}. Think of the SSE as a landscape
over all possible knob settings, with hills (bad fits) and valleys (good
fits); the optimizer is a hiker looking for the lowest point. Problems
arise when the landscape has a long flat valley (two knobs can
compensate each other), a flat plateau (the hiker cannot tell which way
is down), or several separate valleys (the hiker may stop in the wrong
one). Most of this week is about recognizing those situations and
protecting yourself against them.

\textbf{Suggested timing (60 min):} discrepancy functions 10 · error
surfaces 10 · optimization algorithms 20 · practical problems and a
fitting checklist 15 · preview 5.

\subsubsection{1. Discrepancy functions}

For observations \(y_i\) and model predictions
\(\hat y_i = f(x_i; \boldsymbol\theta)\), \(i = 1, \dots, N\):

\[
\mathrm{SSE}(\boldsymbol\theta) = \sum_{i=1}^{N} \big(y_i - f(x_i;\boldsymbol\theta)\big)^2, \qquad
\mathrm{MSE} = \frac{\mathrm{SSE}}{N}, \qquad
\mathrm{RMSD} = \sqrt{\mathrm{MSE}}.
\]

\begin{itemize}
\tightlist
\item
  All three have the \textbf{same minimizer}: they are monotone
  transformations of each other. We \emph{minimize} SSE and
  \emph{report} RMSD, because RMSD is in the units of the data (here:
  proportion recalled), so ``RMSD = 0.057'' means predictions are off by
  about 0.06 on average.
\item
  The \textbf{least-squares estimate} is \[
  \hat{\boldsymbol\theta}_{\text{LS}} = \arg\min_{\boldsymbol\theta} \mathrm{SSE}(\boldsymbol\theta).
  \]
\end{itemize}

\textbf{Why squares?} Squaring makes all errors positive, penalizes
large errors heavily, and gives a smooth (differentiable) function that
optimizers handle well. More deeply, minimizing SSE is exactly maximum
likelihood estimation \textbf{if} the noise is Gaussian with the same
variance for every observation (Week 4). That ``if'' is the hidden
assumption of least squares.

\textbf{When the assumption fails.} Recall proportions are binomial:
their variance \(f(1-f)/n\) depends on \(f\). Points near 0.5 are
noisier than points near 0 or 1, but plain SSE treats them as equally
reliable. \textbf{Weighted least squares} fixes this by weighting each
squared error by the inverse of its variance:

\[
\mathrm{WSSE}(\boldsymbol\theta) = \sum_{i} \frac{\big(y_i - \hat y_i\big)^2}{\mathrm{Var}(y_i)}.
\]

You will try this in Exercise 3. Next week, maximum likelihood handles
this issue more naturally.

\begin{quote}
\textbf{Note:} some textbooks call RMSD ``RMSE''. They are the same
quantity.
\end{quote}

\paragraph{Reading the math}

\begin{itemize}
\tightlist
\item
  \(\sum_{i=1}^{N}\) means ``add up over all \(N\) data points''.
  \(y_i - f(x_i;\boldsymbol\theta)\) is the \textbf{residual} of point
  \(i\): observed minus predicted. Squaring and adding gives one number
  per parameter setting \(\boldsymbol\theta\).
\item
  \(\arg\min_{\boldsymbol\theta} \mathrm{SSE}(\boldsymbol\theta)\) reads
  ``the value of \(\boldsymbol\theta\) at which SSE is smallest''. Note
  the difference: \(\min\) is the smallest \emph{SSE value},
  \(\arg\min\) is the \emph{parameter value} that achieves it. We want
  the \(\arg\min\).
\item
  \textbf{Small example.} Data \(y = (0.62, 0.41)\), predictions
  \((0.66, 0.44)\): residuals \((-0.04, -0.03)\), SSE
  \(= 0.0016 + 0.0009 = 0.0025\), MSE \(= 0.00125\), RMSD \(= 0.035\):
  the curve misses the points by about 3.5 percentage points on average.
\item
  In WSSE each squared residual is divided by the variance of that
  point, so a noisy point counts less. With binomial data the variance
  is \(f(1-f)/n\).
\end{itemize}

\subsubsection{2. The error surface}

For a model with \(k\) parameters, \(\mathrm{SSE}(\boldsymbol\theta)\)
is a surface over a \(k\)-dimensional parameter space. With \(k = 2\) we
can draw it as a contour map: the fitting problem becomes ``find the
lowest point of the landscape.''

Features to recognize:

\begin{longtable}[]{@{}
  >{\raggedright\arraybackslash}p{(\linewidth - 4\tabcolsep) * \real{0.3333}}
  >{\raggedright\arraybackslash}p{(\linewidth - 4\tabcolsep) * \real{0.3333}}
  >{\raggedright\arraybackslash}p{(\linewidth - 4\tabcolsep) * \real{0.3333}}@{}}
\caption{\centering 2. The error surface}\tabularnewline
\toprule\noalign{}
\begin{minipage}[b]{\linewidth}\raggedright
Feature
\end{minipage} & \begin{minipage}[b]{\linewidth}\raggedright
What it looks like
\end{minipage} & \begin{minipage}[b]{\linewidth}\raggedright
What it means
\end{minipage} \\
\midrule\noalign{}
\endfirsthead
\toprule\noalign{}
\begin{minipage}[b]{\linewidth}\raggedright
Feature
\end{minipage} & \begin{minipage}[b]{\linewidth}\raggedright
What it looks like
\end{minipage} & \begin{minipage}[b]{\linewidth}\raggedright
What it means
\end{minipage} \\
\midrule\noalign{}
\endhead
\bottomrule\noalign{}
\endlastfoot
\textbf{Single, round bowl} & Concentric, nearly circular contours &
Parameters well determined and roughly independent; any decent algorithm
finds the minimum \\
\textbf{Long, diagonal valley} & Elongated contours tilted at an angle &
\textbf{Parameter trade-off}: increasing one parameter can be
compensated by changing the other. Estimates will be correlated and
uncertain along the valley (Week 5, Week 7) \\
\textbf{Plateau} & Very widely spaced contours & The model's predictions
barely change with the parameters (e.g., a decay rate so large that the
predicted curve is 0 everywhere). Gradient-based methods stall \\
\textbf{Several separate basins} & More than one closed set of contours
& \textbf{Local minima}: the algorithm's answer depends on where it
starts \\
\end{longtable}

For exponential forgetting, \(A\) and \(\lambda\) trade off: a higher
starting level with faster decay can predict nearly the same mid-range
values as a lower level with slower decay. The valley is therefore
tilted (you will see it in the lab).

\subsubsection{3. Optimization algorithms}

All algorithms below start from initial values
\(\boldsymbol\theta^{(0)}\) and generate a sequence
\(\boldsymbol\theta^{(1)}, \boldsymbol\theta^{(2)}, \dots\) that
(hopefully) descends to the minimum.

\paragraph{3.1 Grid search}

Evaluate the SSE at every point of a regular grid and keep the best.

\begin{itemize}
\tightlist
\item
  \textbf{Pros:} trivially simple; cannot be fooled by local minima
  \emph{on the grid}; gives a picture of the whole surface.
\item
  \textbf{Cons:} precision limited by grid spacing; cost grows
  \textbf{exponentially} with the number of parameters. With 100 values
  per parameter: \(k = 2\) needs \(10^4\) evaluations, \(k = 5\) needs
  \(10^{10}\). This is the \textbf{curse of dimensionality}.
\item
  \textbf{Use it for:} visualizing 1--2-parameter surfaces, and for
  generating good starting values.
\end{itemize}

\paragraph{3.2 Nelder--Mead simplex (derivative-free)}

A \textbf{simplex} is a set of \(k+1\) points in a \(k\)-dimensional
space (a triangle when \(k = 2\)). Each iteration tries to replace the
\textbf{worst} vertex \(\boldsymbol\theta_w\) by moving it through the
centroid \(\bar{\boldsymbol\theta}\) of the other vertices:

\begin{longtable}[]{@{}
  >{\raggedright\arraybackslash}p{(\linewidth - 4\tabcolsep) * \real{0.3333}}
  >{\raggedright\arraybackslash}p{(\linewidth - 4\tabcolsep) * \real{0.3333}}
  >{\raggedright\arraybackslash}p{(\linewidth - 4\tabcolsep) * \real{0.3333}}@{}}
\caption{\centering 3.2 Nelder--Mead simplex (derivative-free)}\tabularnewline
\toprule\noalign{}
\begin{minipage}[b]{\linewidth}\raggedright
Move
\end{minipage} & \begin{minipage}[b]{\linewidth}\raggedright
Formula (standard coefficients)
\end{minipage} & \begin{minipage}[b]{\linewidth}\raggedright
When
\end{minipage} \\
\midrule\noalign{}
\endfirsthead
\toprule\noalign{}
\begin{minipage}[b]{\linewidth}\raggedright
Move
\end{minipage} & \begin{minipage}[b]{\linewidth}\raggedright
Formula (standard coefficients)
\end{minipage} & \begin{minipage}[b]{\linewidth}\raggedright
When
\end{minipage} \\
\midrule\noalign{}
\endhead
\bottomrule\noalign{}
\endlastfoot
Reflect &
\(\boldsymbol\theta_r = \bar{\boldsymbol\theta} + (\bar{\boldsymbol\theta} - \boldsymbol\theta_w)\)
& Always tried first \\
Expand &
\(\boldsymbol\theta_e = \bar{\boldsymbol\theta} + 2(\bar{\boldsymbol\theta} - \boldsymbol\theta_w)\)
& The reflected point is better than the current best; try going
further \\
Contract &
\(\boldsymbol\theta_c = \bar{\boldsymbol\theta} + \tfrac{1}{2}(\boldsymbol\theta_w - \bar{\boldsymbol\theta})\)
(inside), or the same step on the reflected side (outside) & The
reflected point is not better than the second-worst vertex \\
Shrink & Move all vertices halfway toward the best one & Contraction
also failed \\
\end{longtable}

The reflected point is \textbf{accepted} when it is better than the
second-worst vertex but not better than the best. The simplex
``tumbles'' downhill, stretches along valleys, and shrinks around the
minimum.

\begin{itemize}
\tightlist
\item
  \textbf{Pros:} needs only SSE values (no derivatives); robust to noisy
  or non-smooth objective functions; the default workhorse in cognitive
  modeling.
\item
  \textbf{Cons:} can be slow in high dimensions; can stall in long
  narrow valleys; no guarantee of finding the global minimum.
\end{itemize}

\begin{quote}
Direct-search methods like Nelder--Mead are treated in depth in
Ghodousian's book on direct methods in nonlinear optimization (see
Further reading).
\end{quote}

\paragraph{3.3 Gradient-based methods}

The \textbf{gradient} \(\nabla \mathrm{SSE}(\boldsymbol\theta)\) points
uphill, so we step the other way:

\[
\boldsymbol\theta^{(m+1)} = \boldsymbol\theta^{(m)} - \eta\,\nabla \mathrm{SSE}\big(\boldsymbol\theta^{(m)}\big),
\qquad
\frac{\partial\,\mathrm{SSE}}{\partial \theta} = -2 \sum_{i} \big(y_i - \hat y_i\big)\,\frac{\partial \hat y_i}{\partial \theta},
\]

with step size (learning rate) \(\eta\). For the exponential model,
\(\partial \hat y/\partial A = e^{-\lambda t}\) and
\(\partial \hat y/\partial \lambda = -A\,t\,e^{-\lambda t}\).

In practice we use \textbf{quasi-Newton} methods such as \textbf{BFGS}.
They also estimate the \emph{curvature} of the surface and choose the
step size automatically, so they converge much faster near a minimum.
\textbf{L-BFGS-B} is the variant that accepts \textbf{bounds} on
parameters.

\begin{itemize}
\tightlist
\item
  \textbf{Pros:} fast and precise on smooth surfaces; scale well to many
  parameters (this is how neural networks are trained).
\item
  \textbf{Cons:} need a smooth objective; get stuck where the gradient
  is (nearly) zero, at plateaus as well as at minima; sensitive to
  starting values.
\end{itemize}

\begin{quote}
\textbf{Advanced: why plateaus trap gradient methods.} If \(\lambda\) is
very large, \(e^{-\lambda t} \approx 0\) at every observed \(t\), so
both partial derivatives above are \(\approx 0\). A gradient method sees
a ``flat floor'' and stops, sometimes reporting successful convergence.
Nelder--Mead compares function values at distant points and can step off
the plateau. You will see both behaviors in lab Part D.
\end{quote}

\textbf{Summary:}

\begin{longtable}[]{@{}
  >{\raggedright\arraybackslash}p{(\linewidth - 6\tabcolsep) * \real{0.2500}}
  >{\raggedright\arraybackslash}p{(\linewidth - 6\tabcolsep) * \real{0.2500}}
  >{\raggedright\arraybackslash}p{(\linewidth - 6\tabcolsep) * \real{0.2500}}
  >{\raggedright\arraybackslash}p{(\linewidth - 6\tabcolsep) * \real{0.2500}}@{}}
\caption{\centering 3.3 Gradient-based methods}\tabularnewline
\toprule\noalign{}
\begin{minipage}[b]{\linewidth}\raggedright
Method
\end{minipage} & \begin{minipage}[b]{\linewidth}\raggedright
Needs derivatives
\end{minipage} & \begin{minipage}[b]{\linewidth}\raggedright
Bounds
\end{minipage} & \begin{minipage}[b]{\linewidth}\raggedright
Typical use
\end{minipage} \\
\midrule\noalign{}
\endfirsthead
\toprule\noalign{}
\begin{minipage}[b]{\linewidth}\raggedright
Method
\end{minipage} & \begin{minipage}[b]{\linewidth}\raggedright
Needs derivatives
\end{minipage} & \begin{minipage}[b]{\linewidth}\raggedright
Bounds
\end{minipage} & \begin{minipage}[b]{\linewidth}\raggedright
Typical use
\end{minipage} \\
\midrule\noalign{}
\endhead
\bottomrule\noalign{}
\endlastfoot
Grid search & No & Built in & Visualization, starting values,
\(k \le 2\) \\
Nelder--Mead & No & Yes ( \ensuremath{\geq} 1.7 & Default for small cognitive models;
noisy objectives \\
BFGS & Numerical or analytic & No & Smooth objectives, unconstrained
parameters \\
L-BFGS-B & Numerical or analytic & Yes & Smooth objectives with bounded
parameters \\
\end{longtable}

\paragraph{Reading the math}

The gradient step
\(\boldsymbol\theta^{(m+1)} = \boldsymbol\theta^{(m)} - \eta\,\nabla \mathrm{SSE}(\boldsymbol\theta^{(m)})\)
reads: \emph{new guess = old guess − step size × slope}. The superscript
\((m)\) is the iteration number, not a power. \(\nabla \mathrm{SSE}\) is
the vector of partial derivatives (one slope per parameter). The minus
sign makes the step go downhill.

The derivative formula
\(\partial\,\mathrm{SSE}/\partial\theta = -2\sum_i (y_i - \hat y_i)\,\partial \hat y_i/\partial\theta\)
follows from the chain rule. It says: points where the residual is large
\emph{and} where the prediction is sensitive to \(\theta\) pull hardest
on \(\theta\). At the minimum these pulls cancel and the sum is zero.

\paragraph{For an engineer}

\begin{itemize}
\tightlist
\item
  \textbf{Least squares} is the same criterion as in channel estimation
  and curve fitting; the difference here is that the models are
  \textbf{nonlinear} in the parameters (\(e^{-\lambda t}\)), so there is
  no closed-form normal-equation solution and we need iterative search.
\item
  \textbf{Weighted least squares} is noise whitening: dividing by the
  variance is what you do before LS when the noise is not white
  (heteroscedastic noise).
\item
  \textbf{Gradient descent with step size \(\eta\)} is the
  steepest-descent idea behind the LMS adaptive filter. BFGS is a
  quasi-Newton method (it builds an approximation of the Hessian), close
  to Gauss--Newton / Levenberg--Marquardt used in nonlinear LS.
\item
  \textbf{A long diagonal valley} in the error surface is an
  \textbf{ill-conditioned} problem: the Hessian has a large eigenvalue
  ratio, as with strongly correlated regressors. Parameters are well
  determined along one direction and poorly along the other.
\item
  \textbf{Reparameterization} (\(\lambda = e^{\eta}\)) is a change of
  variables that turns a constrained problem into an unconstrained one;
  it also often improves conditioning.
\item
  \textbf{What is different from engineering practice:} data sets are
  small (8 points, 50 items each), the noise is large, and the
  parameters must be \emph{interpreted} psychologically. So it is not
  enough that the predicted curve looks good: the individual parameter
  values must be trustworthy.
\end{itemize}

\subsubsection{4. Practical problems and how to handle them}

\textbf{1. Starting values.} Use informed starts: read them off a plot,
use a quick linearized fit (Week 2), or take the best point of a coarse
grid.

\textbf{2. Local minima \ensuremath{\rightarrow} multi-start.} Run the optimizer from many
random starting points and keep the best solution. If many starts agree
on the same minimum, you can be reasonably confident it is the global
one. Multi-start is cheap insurance; use it routinely for models with
more than two parameters.

\textbf{3. Parameter bounds.} Many parameters have natural ranges
(\(0 \le A \le 1\), \(\lambda > 0\)). Without constraints, optimizers
happily return \(A = 29.6\). Three approaches:

\begin{longtable}[]{@{}
  >{\raggedright\arraybackslash}p{(\linewidth - 4\tabcolsep) * \real{0.3333}}
  >{\raggedright\arraybackslash}p{(\linewidth - 4\tabcolsep) * \real{0.3333}}
  >{\raggedright\arraybackslash}p{(\linewidth - 4\tabcolsep) * \real{0.3333}}@{}}
\caption{\centering 4. Practical problems and how to handle them}\tabularnewline
\toprule\noalign{}
\begin{minipage}[b]{\linewidth}\raggedright
Approach
\end{minipage} & \begin{minipage}[b]{\linewidth}\raggedright
How
\end{minipage} & \begin{minipage}[b]{\linewidth}\raggedright
Comment
\end{minipage} \\
\midrule\noalign{}
\endfirsthead
\toprule\noalign{}
\begin{minipage}[b]{\linewidth}\raggedright
Approach
\end{minipage} & \begin{minipage}[b]{\linewidth}\raggedright
How
\end{minipage} & \begin{minipage}[b]{\linewidth}\raggedright
Comment
\end{minipage} \\
\midrule\noalign{}
\endhead
\bottomrule\noalign{}
\endlastfoot
Bounded optimizer & or with & Easiest; recommended default \\
\textbf{Reparameterization} & Optimize over unconstrained \(\eta\) and
transform: \(\lambda = e^{\eta}\) (positive),
\(A = 1/(1 + e^{-\alpha_A})\) (between 0 and 1) & Works with any
optimizer; also used in Bayesian modeling (Week 13) \\
Penalty & Return a huge SSE when parameters are out of range & Creates
discontinuities; avoid \\
\end{longtable}

\textbf{4. Non-identifiability.} Sometimes \emph{several} parameter
combinations fit equally well. An example is a two-component model whose
two components can swap labels. This is a property of the model, not a
failure of the optimizer, and multi-start reveals it (lab Part D).

Always (i) plot the fitted curve over the data, (ii) look at the
residuals, and (iii) check that the estimates are within plausible
ranges.

\textbf{A fitting checklist (use it in every lab from now on):}

\begin{enumerate}
\def\labelenumi{\arabic{enumi}.}
\tightlist
\item
  Plot the data.
\item
  Write the model and the discrepancy function; test them on a few
  hand-picked parameter values.
\item
  Choose sensible bounds and starting values.
\item
  Optimize, preferably with multi-start.
\item
  Check the convergence message, the plot, the residuals and the
  parameter ranges.
\item
  Run a recovery check: simulate data from the fitted parameters, refit,
  and compare (Week 7 makes this systematic).
\end{enumerate}

\begin{quote}
\textbf{Advanced: conditionally linear parameters.} In
\(f(t) = A\,e^{-\lambda t}\), for any fixed \(\lambda\) the best \(A\)
has a closed form, because the model is linear in \(A\):
\(\hat A(\lambda) = \sum_i y_i g_i / \sum_i g_i^2\) with
\(g_i = e^{-\lambda t_i}\). Substituting it back gives a one-dimensional
problem in \(\lambda\) alone. This trick (``profiling out'' linear
parameters, or \emph{variable projection}) makes fitting faster and more
stable, and the resulting 1-D function is related to the \emph{profile
likelihood} of Week 5.
\end{quote}

\paragraph{For the instructor}

\begin{itemize}
\tightlist
\item
  \textbf{Show the landscape first.} A contour plot of SSE over
  \((A, \lambda)\) (lab Part B) makes every later idea concrete: the
  minimum, the diagonal valley, the plateau. Put it on the projector
  early and keep coming back to it.
\item
  \textbf{Act out Nelder--Mead.} Draw a triangle on the contour plot and
  do one reflection by hand (worked example Step 2). Students remember
  the ``tumbling triangle'' much better than the table of moves.
\item
  \textbf{Your strength:} you know gradient descent and least squares.
  Spend less time on the algorithms and more on the \emph{failure modes}
  (Part D of the lab): an optimizer reporting success at \(A = 29.6\) is
  the most memorable lesson of the week.
\item
  \textbf{Typical questions.}

  \begin{itemize}
  \tightlist
  \item
    \emph{``Which optimizer should I use?''} For small cognitive models:
    Nelder--Mead or L-BFGS-B with bounds, always with multi-start.
  \item
    \emph{``Why not always use a grid?''} The cost grows exponentially
    with the number of parameters.
  \item
    \emph{``Is SSE the right measure?''} Only if noise is Gaussian with
    equal variance; next week replaces it with the likelihood.
  \end{itemize}
\end{itemize}

\begin{center}\rule{0.5\linewidth}{0.5pt}\end{center}

\subsection{Worked example}

\textbf{Data.} A participant recalls proportions
\(y = (0.62, 0.41, 0.19)\) at \(t = (1, 3, 7)\) days. Model:
\(f(t) = A\,e^{-\lambda t}\).

\subsubsection{Step 1: a small grid search by hand}

For \((A, \lambda) = (0.8, 0.2)\): predictions \(0.8e^{-0.2} = 0.655\),
\(0.8e^{-0.6} = 0.439\), \(0.8e^{-1.4} = 0.197\), so

\[
\mathrm{SSE} = (0.62 - 0.655)^2 + (0.41 - 0.439)^2 + (0.19 - 0.197)^2 = 0.00123 + 0.00084 + 0.00005 \approx 0.00212.
\]

Repeating this for a 2 × 3 grid:

\begin{longtable}[]{@{}llll@{}}
\caption{\centering Step 1: a small grid search by hand}\tabularnewline
\toprule\noalign{}
& \(\lambda = 0.15\) & \(\lambda = 0.20\) & \(\lambda = 0.25\) \\
\midrule\noalign{}
\endfirsthead
\toprule\noalign{}
& \(\lambda = 0.15\) & \(\lambda = 0.20\) & \(\lambda = 0.25\) \\
\midrule\noalign{}
\endhead
\bottomrule\noalign{}
\endlastfoot
\(A = 0.8\) & 0.02281 & \textbf{0.00212} & 0.00364 \\
\(A = 0.9\) & 0.06638 & 0.02172 & 0.00791 \\
\end{longtable}

The grid's best point is \((0.8, 0.20)\). Note the tilt: \((0.9, 0.25)\)
beats \((0.9, 0.15)\) by far. A larger \(A\) ``wants'' a larger
\(\lambda\), which is the trade-off valley.

\subsubsection{Step 2: one Nelder--Mead iteration by hand}

Start with the simplex (triangle) of three grid points:

\begin{longtable}[]{@{}llll@{}}
\caption{\centering Step 2: one Nelder--Mead iteration by hand}\tabularnewline
\toprule\noalign{}
Vertex & \((A, \lambda)\) & SSE & Rank \\
\midrule\noalign{}
\endfirsthead
\toprule\noalign{}
Vertex & \((A, \lambda)\) & SSE & Rank \\
\midrule\noalign{}
\endhead
\bottomrule\noalign{}
\endlastfoot
\(\boldsymbol\theta_b\) & (0.8, 0.25) & 0.00364 & best \\
\(\boldsymbol\theta_g\) & (0.8, 0.15) & 0.02281 & second-worst \\
\(\boldsymbol\theta_w\) & (0.9, 0.15) & 0.06638 & worst \\
\end{longtable}

\begin{enumerate}
\def\labelenumi{\arabic{enumi}.}
\tightlist
\item
  \textbf{Centroid} of all vertices except the worst:
  \(\bar{\boldsymbol\theta} = \big((0.8 + 0.8)/2,\ (0.25 + 0.15)/2\big) = (0.8, 0.20)\).
\item
  \textbf{Reflect} the worst vertex through the centroid:
  \(\boldsymbol\theta_r = (0.8, 0.20) + \big((0.8, 0.20) - (0.9, 0.15)\big) = (0.7, 0.25)\).
\item
  \textbf{Evaluate:} predictions \((0.545, 0.331, 0.122)\), so
  \(\mathrm{SSE}(\boldsymbol\theta_r) = 0.01657\).
\item
  \textbf{Decide:} \(0.00364 \le 0.01657 < 0.02281\). The reflected
  point is better than the second-worst but not better than the best, so
  we \textbf{accept the reflection} (no expansion).
\item
  \textbf{New simplex:} \(\{(0.8, 0.25),\ (0.7, 0.25),\ (0.8, 0.15)\}\),
  and the worst SSE has dropped from 0.066 to 0.023.
\end{enumerate}

\subsubsection{Step 3: let the computer finish}

\(\hat A = 0.754\) · \(\hat\lambda = 0.200\) ·
\(\mathrm{SSE} = 3.7 \times 10^{-5}\) · \(= 0.0035\)

\textbf{Interpretation.} The fitted model says this participant would
recall about 75\% immediately after learning and loses about 18\% of the
remaining material per day (\(1 - e^{-0.2} = 0.181\)), with a half-life
of \(\ln 2 / 0.2 \approx 3.5\) days. With only three data points and two
parameters, the near-perfect fit says very little about whether the
exponential model is \emph{right}. Any reasonable two-parameter
decreasing curve would fit three points well (Week 6).

\begin{center}\rule{0.5\linewidth}{0.5pt}\end{center}

\Needspace{11\baselineskip}

\subsection{Lab (2 hours)}

\subsubsection{Goals}

\begin{enumerate}
\def\labelenumi{\arabic{enumi}.}
\tightlist
\item
  Write an SSE function and explore the error surface of two forgetting
  models.
\item
  Fit by grid search, then with Nelder--Mead and L-BFGS-B; compare with
  last week's linearized fits.
\item
  Provoke and fix fitting failures: plateaus, local minima and
  out-of-range estimates.
\item
  \textbf{Evaluate} the fits (RMSD, residuals, comparison with the true
  values) and \textbf{interpret} the parameters.
\end{enumerate}

\textbf{Timing:} Part A 15 min · Part B 25 min · Part C 25 min · Part D
35 min · Part E 20 min.

\subsubsection{Part A: Get a feel for the discrepancy function (15 min)}

Before running the next parts, \textbf{predict}: which parameter values
will give the smallest SSE for each model? Write your guess down.

\subsubsection{Part B: Grid search and the error surface (25 min)}

\textbf{What to notice:} both surfaces have a single basin, but it is an
\textbf{elongated, tilted valley}: higher \(A\) goes with higher decay.
That is the parameter trade-off. For the power model the true value
(white ×) lies inside the valley but not at its bottom: with finite
data, the best fit is never exactly the truth.

\subsubsection{Part C: Fitting with an optimizer (25 min)}

\textbf{What to notice:} - Nelder--Mead and L-BFGS-B agree to three or
four decimals; the grid is within its spacing (0.005 for \(A\)) of them.
The optimizers needed only 36--71 function evaluations instead of 42
300. - Least squares gives a \textbf{lower SSE} than the linearized fits
of Week 2 (it is minimizing exactly that quantity), and the power
estimates are much closer to the truth. - The path first heads quickly
into the valley and then crawls along it, a typical pattern.

\subsubsection{Part D: When fitting goes wrong (35 min)}

\textbf{D1. A plateau and a false ``success''.} Start the exponential
fit with an absurd decay rate, \(\lambda = 20\) per day:

They started on a plateau where the gradient is almost zero (Lecture
§3.3, Advanced box). Nelder--Mead escaped.

\textbf{D2. Reparameterization.} Make the parameters unconstrained so
that \emph{any} optimizer respects the ranges \(0 < A < 1\) and
\(\lambda > 0\):

Same answer as in Part C, and the estimates are guaranteed to lie in
range.

\textbf{D3. Local minima and multi-start.} A popular idea in memory
research is that forgetting reflects two memory traces: a \textbf{fast}
component that decays quickly and a \textbf{slow} one that lasts:

\[
f(t) = A_1 e^{-\lambda_1 t} + A_2 e^{-\lambda_2 t}.
\]

This model has \(k = 4\) parameters, and its error surface has several
basins. Fit it from 50 random starting points:

\textbf{What to notice:} - About 40 of 50 starts end in the same best
basin, but the rest (about 1 in 6) get stuck in \textbf{worse local
minima}, some with an SSE many times larger. - Among the best solutions,
the fast and slow components may appear \textbf{with labels swapped}
(\(A_1, \lambda_1 \leftrightarrow A_2, \lambda_2\)). Both describe the
same curve. This is a non-identifiability built into the model, and a
common fix is to impose an order constraint such as
\(\lambda_1 > \lambda_2\). - In the best solution one component has a
rate of exactly 0 (a flat, non-decaying component of 0.264). The fitted
two-trace model has therefore collapsed into the
exponential-with-asymptote model of Exercise 4, with the same SSE
(0.0201). Is that a meaningful ``permanent memory'' or just the extra
parameters absorbing noise? Remember that the data came from a
\emph{power} model. Fit alone cannot tell (Weeks 6--7).

\subsubsection{Part E: Evaluate and interpret (20 min)}

\begin{enumerate}
\def\labelenumi{\arabic{enumi}.}
\tightlist
\item
  Compare the residual patterns. Which model shows a systematic pattern,
  and what does that pattern say about the shape of forgetting in these
  data?
\item
  The power model's relative loss rate falls from about 0.21 per day
  after 1 day to about 0.02 per day after 20 days. Explain in plain
  language what this claims about memory.
\item
  The least-squares estimates (\(\hat A \approx 0.87\),
  \(\hat d \approx 0.43\)) are closer to the truth than last week's
  linearized estimates but still not equal to it. Is this a failure of
  the method? (Hint: what would happen if you reran the experiment with
  a new random sample?)
\end{enumerate}

\subsubsection{Exercises (increasing difficulty)}

\textbf{Exercise 1 \ensuremath{\star} (same minimizer).} Fit the power model by
minimizing SSE, MSE and RMSD in turn. Show that the estimates agree (up
to optimizer tolerance). Why do we still prefer to \emph{report} RMSD?

\textbf{Exercise 2 \ensuremath{\star} (Ebbinghaus by least squares).} Refit the
Ebbinghaus savings data from Week 2, Exercise 3 (values marked
{[}verify{]}) with least squares, for both models. Compare the
parameters and RMSD with your linearized fits from Week 2.

\textbf{Exercise 3 \ensuremath{\star}\ensuremath{\star} (weighted least squares).} Binomial proportions
have variance \(f(1-f)/n\). Formulate and evaluate WSSE (Lecture §1) using the
\textbf{predicted} \(f\) in the variance, refit the power model and
compare the estimates with plain least squares. Which data points gain
influence and which lose it? (Next week you will see that maximum
likelihood does this weighting automatically.)

\textbf{Exercise 4 \ensuremath{\star}\ensuremath{\star} (more parameters, better fit?).} Fit the
three-parameter \textbf{exponential-with-asymptote} model
\(f(t) = B + (A - B)e^{-\lambda t}\), where \(B\) is the long-term
asymptote, using bounds \(0 \le B \le A \le 1\) (hint: reparameterize,
or check the constraint after fitting). \emph{Expected:} SSE \ensuremath{\approx} 0.020,
\textbf{lower} than the power model's \ensuremath{\approx} 0.026, even though the data were
generated by the power model! What does this tell you about using SSE to
choose between models with different numbers of parameters? (We answer
this properly in Weeks 6--7.)

\textbf{Exercise 5 \ensuremath{\star}\ensuremath{\star}\ensuremath{\star} (Advanced: Nelder--Mead from scratch).} Formulate and evaluate
Nelder--Mead for \(k = 2\) using the moves in Lecture §3.2 (reflect,
expand, inside/outside contract, shrink) and a stopping rule based on
the spread of SSE values in the simplex.

\subsubsection{Lab check-off}

Show the instructor or TA: (1) the error-surface plots with grid minima
(Part B), (2) the comparison table from Part C, (3) your multi-start
histogram (Part D3) with one sentence on what it shows, and (4) your
answer to interpretation question 1 (Part E).

\begin{center}\rule{0.5\linewidth}{0.5pt}\end{center}

\subsection{Connection to AI \& LLMs}

\textbf{Scaling laws are curve fitting.} One of the most influential
results in modern AI is that the test loss of language models falls
smoothly with model size \(N\), data size and compute, following a
\textbf{power law} with an asymptote: \(L(N) = E + B\,N^{-\alpha}\)
(Kaplan et al., 2020; Hoffmann et al., 2022). Companies use such fits to
decide how big the next model should be. Mathematically, this is the
same problem as fitting a forgetting curve: a nonlinear model with three
parameters, fitted by least squares with a numerical optimizer and
multiple starting points. It has the same pitfalls too: extrapolation
far beyond the data, and trade-offs between parameters (here \(E\) and
\(\alpha\)).

\(E = 0\) Compare the SSE and the extrapolation to \(N = 10^{11}\).
Which conclusion about ``bigger is better'' changes? - \ensuremath{\star}\ensuremath{\star} Fit only the
four smallest models and extrapolate to the four largest. How far off is
the prediction? Repeat with different seeds to see how stable the
extrapolation is. - \ensuremath{\star}\ensuremath{\star} Plot the SSE surface over \(E\) and \(\alpha\)
(with \(B\) at its best value for each pair). Is there a ridge, as in
the Week 3 lab?

\emph{Optional reading:} Kaplan, J., et al.~(2020). Scaling laws for
neural language models. \emph{arXiv:2001.08361}. Hoffmann, J., et
al.~(2022). Training compute-optimal large language models.
\emph{arXiv:2203.15556}.

\begin{center}\rule{0.5\linewidth}{0.5pt}\end{center}

\subsection{Common mistakes \& misconceptions}

\begin{enumerate}
\def\labelenumi{\arabic{enumi}.}
\tightlist
\item
  Always inspect estimates, plots and residuals.
\item
  \textbf{One start is enough.} For anything beyond simple two-parameter
  models, run multi-start. The cost is small; the risk of reporting a
  local minimum is not.
\item
  \textbf{Ignoring parameter ranges.} Unbounded optimizers can return
  impossible values (negative rates, probabilities above 1). Use bounds
  or reparameterization.
\item
  \textbf{Comparing SSE across models with different numbers of
  parameters.} The more flexible model almost always wins on SSE
  (Exercise 4). That is a question of model \emph{evaluation}, not
  estimation (Weeks 6--7).
\item
  \textbf{Comparing SSE across datasets of different sizes.} SSE grows
  with \(N\); use RMSD for comparisons.
\item
  \textbf{Treating the minimizer as the true value.}
  \(\hat{\boldsymbol\theta}\) is an estimate with sampling error. A
  different sample would give a different estimate (Week 5 quantifies
  this).
\item
  \textbf{Forgetting what least squares assumes.} Plain SSE assumes
  equal, Gaussian noise for all observations. Proportions, counts and
  RTs usually violate this.
\item
  \textbf{Using grid search for many parameters.} The cost explodes
  (\(100^k\) evaluations). Use grids only for visualization and starting
  values.
\end{enumerate}

\begin{center}\rule{0.5\linewidth}{0.5pt}\end{center}

\subsection{Check your understanding}

\textbf{1.} Why do SSE, MSE and RMSD have the same minimizer? Which one
would you report, and why? \textbf{Answer} MSE = SSE/N and RMSD = \ensuremath{\surd}MSE
are strictly increasing transformations of SSE (N is fixed), so the same
parameters minimize all three. Report RMSD, because it is in the units
of the data and comparable across datasets of different sizes.

\textbf{2.} What is a parameter trade-off, and how does it appear on an
error surface? \textbf{Answer} Changes in one parameter can be
compensated by changes in another, leaving the predictions (and the SSE)
almost unchanged. On the surface it appears as an elongated, tilted
valley of low SSE. Estimates are then correlated and poorly determined
along the valley direction.

\textbf{3.} With 50 grid values per parameter, how many SSE evaluations
does a grid search need for a model with 6 parameters? What does this
imply? \textbf{Answer} \(50^6 \approx 1.56 \times 10^{10}\) evaluations.
Grid search is infeasible beyond 2--3 parameters (the curse of
dimensionality); we need iterative optimizers.

\textbf{4.} In one Nelder--Mead iteration, the reflected point is better
than the current best vertex. What does the algorithm do next?
\textbf{Answer} It tries an expansion: a point twice as far from the
centroid in the same direction. It keeps the expanded point if it is
better than the reflected point, otherwise the reflected point.

\textbf{5.} Why can a gradient-based optimizer stop on a plateau and
report success? \textbf{Answer} Its stopping rule is based on the
gradient (or step size) being nearly zero. On a plateau the gradient is
close to zero even though the point is not a minimum, so the convergence
criterion is met.

\textbf{6.} Give two ways to keep a decay-rate parameter positive during
fitting. (ii) Reparameterize: optimize \(\eta\) on the whole real line
and set \(\lambda = e^{\eta}\).

\textbf{7.} Multi-start fits of the two-component model reach the same
minimum SSE with
\((A_1, \lambda_1, A_2, \lambda_2) = (0.58, 0.36, 0.26, 0.00)\) and
\((0.26, 0.00, 0.58, 0.36)\). Is this a problem with the optimizer?
\textbf{Answer} No.~It is label switching: swapping the two components
gives the identical curve, so the model is not identifiable without an
ordering constraint (e.g., \(\lambda_1 > \lambda_2\)). The optimizer
correctly found two equivalent global minima.

\begin{center}\rule{0.5\linewidth}{0.5pt}\end{center}

\subsection{Assignment}

No new graded assignment is released this week. \textbf{A1 (From verbal
to formal) is due at the start of the Week 4 lecture.} Today's lab gives
you everything needed for A1's calibration step if you want to check
your closed-form derivation numerically.

\textbf{Optional practice (ungraded, recommended before Week 4):}
complete lab Exercises 3 and 4. Exercise 3 (weighted least squares) is
the bridge to maximum likelihood, and Exercise 4 previews the central
problem of model evaluation. Solutions will be discussed briefly at the
start of the Week 4 lab.

\begin{center}\rule{0.5\linewidth}{0.5pt}\end{center}

\subsection{Further reading}

\begin{itemize}
\tightlist
\item
  \textbf{Farrell \& Lewandowsky (2018), chapter 3.} Basic parameter
  estimation: discrepancy functions, error surfaces, the simplex
  algorithm, and practical issues such as local minima and starting
  values. \emph{(Main reading.)}
\item
  \textbf{Ghodousian (1401)}, \emph{Mabani-ye ravesh-haye mostaghim dar
  behine-sazi gheyr-khatti} (in Persian): direct search methods for
  nonlinear optimization, including simplex-type methods. Recommended
  for students who want the theory behind Nelder--Mead {[}verify
  chapter{]}.
\item
  \textbf{Verguts (2022)}: the chapter on optimization and gradient
  descent, useful for the gradient-based view that returns in Week 15
  (learning models) {[}verify chapter{]}.
\end{itemize}

\emph{Preview for Week 4:} instead of asking ``which parameters minimize
squared error?'', we ask ``which parameters make the observed data most
probable?'' This is maximum likelihood. It handles binomial data
correctly, and it is the foundation for AIC/BIC (Week 7) and Bayesian
inference (Weeks 9--13).
